\documentclass[11pt]{article}

\usepackage[margin=0.8in]{geometry}
\usepackage{graphicx}
\usepackage{booktabs}
\usepackage{array}
\usepackage{float}
\usepackage{listings}
\usepackage{xcolor}
\usepackage{hyperref}
\usepackage{setspace}

\hypersetup{colorlinks=true, linkcolor=blue, urlcolor=blue}
\lstdefinestyle{code}{
  basicstyle=\ttfamily\small,
  breaklines=true,
  frame=single,
  columns=fullflexible,
  keepspaces=true,
  showstringspaces=false,
  backgroundcolor=\color{gray!8},
  rulecolor=\color{gray!40}
}
\setlength{\parskip}{0pt}
\setlength{\parindent}{0.5in}

\title{CSE 565 Performance Testing Project}
\author{Chase Coleman}
\date{}

\begin{document}
\begin{titlepage}
  \centering
  \doublespacing
  \vspace*{1in}
  {\Large\textbf{CSE 565 Performance Testing Project}\par}
  \vspace{0.5in}
  Chase Coleman\par
  Arizona State University\par
  CSE 565: Software Verification and Validation\par
  Performance Testing Project\par
  \today\par
\end{titlepage}

\setcounter{page}{2}
\doublespacing

\section{Introduction}

This draft documents the performance-testing framework, the three baseline test cases, and the raw execution results. The analysis, conclusions, and improved stress-test strategy will be added after review of the Locust reports.

\section{Task 1: Identification of the Performance Testing Framework}

Locust 2.46.6 was selected as the performance-testing framework. Locust is a Python framework in which a user class defines the work performed by a simulated user. It records request count, failures, response times, response-time percentiles, and requests per second (RPS). The framework can be run from the command line in headless mode and can create HTML and CSV reports.

The official Locust training module used for the framework setup is the ``Your first test'' quick-start guide: \url{https://docs.locust.io/en/stable/quickstart.html}. The guide demonstrates defining a user class in \texttt{locustfile.py}, starting Locust, and reviewing response-time and RPS information. The official headless-mode guide used for the command-line test execution is \url{https://docs.locust.io/en/stable/running-without-web-ui.html}.

The instructor-provided \texttt{load\_test\_tasks.py} file was preserved without modification. The separate \texttt{locustfile.py} defines a Locust user workflow with the same six workload categories: data processing, simulated database query, API request, file I/O, computation, and logging. A controlled local profile was used so that the 10,000-user test could be executed without directing high-volume traffic to the public API in the starter script.

\section{Task 2: Test Case Design}

Each Locust user repeatedly executes the six workload categories in one workflow. Locust records a separate \texttt{TASK} result for each category. The controlled API task returns a simulated response; the remaining operations perform local processing, waiting, file I/O, computation, and logging. The three baseline cases use the same standard profile and a 30-second observation window.

\begin{table}[H]
\centering
\caption{Baseline Locust test cases.}
\begin{tabular}{>{\raggedright\arraybackslash}p{0.85in}>{\raggedright\arraybackslash}p{1.1in}rrrr}
\toprule
\textbf{Test} & \textbf{Load level} & \textbf{Users} & \textbf{Spawn rate} & \textbf{Duration} & \textbf{Profile} \\
\midrule
Test 1 & Baseline load & 100 & 20 users/s & 30 s & Standard \\
Test 2 & Medium load & 1,000 & 200 users/s & 30 s & Standard \\
Test 3 & High load & 10,000 & 1,000 users/s & 30 s & Standard \\
\bottomrule
\end{tabular}
\end{table}

The following command is implemented by \texttt{run\_locust\_scenarios.sh}. It invokes Locust in headless mode and saves an HTML report, statistics CSV, statistics-history CSV, failures CSV, exceptions CSV, and console log for each load level.

\begin{lstlisting}[style=code,language=bash]
CSE565_PROFILE=standard CSE565_API_MODE=simulated
CSE565_DURATION_SECONDS=30 .venv/bin/locust -f locustfile.py --headless
  --users 100 --spawn-rate 20 --stop-timeout 15
  --csv locust_results/users-100/locust
  --html locust_results/users-100/locust_report.html
\end{lstlisting}

\section{Task 3: Test Execution Results}

The three baseline tests completed through Locust on October 4, 2026. Table~\ref{tab:aggregate-results} reproduces the aggregate row from each Locust \texttt{locust\_stats.csv} report. Response-time values are in milliseconds. The result files are retained in the project under \texttt{locust\_results/20261004-123006/}.

\begin{table}[H]
\centering
\caption{Aggregate Locust results for the three baseline tests.}
\label{tab:aggregate-results}
\resizebox{\textwidth}{!}{%
\begin{tabular}{>{\raggedright\arraybackslash}p{0.85in}rrrrrr}
\toprule
\textbf{Users} & \textbf{Requests} & \textbf{Failures} & \textbf{Average (ms)} & \textbf{Median (ms)} & \textbf{Maximum (ms)} & \textbf{RPS} \\
\midrule
100 & 21,512 & 0 & 2.917 & 0.009 & 22.343 & 740.632 \\
1,000 & 221,958 & 0 & 2.705 & 0.009 & 36.228 & 7,390.853 \\
10,000 & 742,403 & 0 & 141.188 & 0.009 & 734.794 & 22,505.491 \\
\bottomrule
\end{tabular}
}
\end{table}

\begin{table}[H]
\centering
\small
\caption{Per-task Locust response-time results.}
\begin{tabular}{>{\raggedright\arraybackslash}p{1.25in}rrr}
\toprule
\textbf{Load and task} & \textbf{Average (ms)} & \textbf{Median (ms)} & \textbf{RPS} \\
\midrule
100: Data Processing & 0.193 & 0.058 & 123.462 \\
100: Database Query & 5.641 & 5.010 & 123.462 \\
100: API Request & 10.916 & 11.000 & 123.427 \\
100: File I/O & 0.554 & 0.082 & 123.427 \\
100: Computation & 0.029 & 0.009 & 123.427 \\
100: Logging & 0.169 & 0.022 & 123.427 \\
\midrule
1,000: Data Processing & 0.082 & 0.057 & 1,232.608 \\
1,000: Database Query & 5.365 & 5.002 & 1,231.809 \\
1,000: API Request & 10.583 & 10.002 & 1,231.609 \\
1,000: File I/O & 0.147 & 0.074 & 1,231.609 \\
1,000: Computation & 0.012 & 0.009 & 1,231.609 \\
1,000: Logging & 0.041 & 0.021 & 1,231.609 \\
\midrule
10,000: Data Processing & 0.063 & 0.057 & 3,812.246 \\
10,000: Database Query & 418.040 & 470.000 & 3,788.025 \\
10,000: API Request & 427.518 & 480.000 & 3,726.305 \\
10,000: File I/O & 0.127 & 0.072 & 3,726.305 \\
10,000: Computation & 0.010 & 0.009 & 3,726.305 \\
10,000: Logging & 0.037 & 0.021 & 3,726.305 \\
\bottomrule
\end{tabular}
\end{table}

\section*{References}

\hangindent=0.5in Locust. (n.d.). \textit{Your first test}. Locust documentation. Retrieved October 4, 2026, from \url{https://docs.locust.io/en/stable/quickstart.html}

\hangindent=0.5in Locust. (n.d.). \textit{Running without the web UI}. Locust documentation. Retrieved October 4, 2026, from \url{https://docs.locust.io/en/stable/running-without-web-ui.html}

\end{document}
